\subsection{nml\_bc}
\label{SEC:nml_bc}
\inp*[nmlBC]{}[nml\_bc][]

This \hyperref[SEC:User_Input]{namelist block}
is used to specify boundary conditions.

\begin{nmlblocklst}[
,escapeinside={!@}{!@}
,label={LST:nml_bc}
,caption={Namelist nml\_bc.}
]
&nml_bc
  Walls_Are_Exact                          = T
  Sub_Inflow_Method                        = 4
  Use_BC_Conflicts                         = T
  BC_Input(1)%Name                         = ""
  BC_Input(1)%Set_BC_Value                 = -9999
  BC_Input(1)%Temperature                  = -1.0               ! [K]
  BC_Input(1)%Pressure                     = -1.0               ! [Pa]
  BC_Input(1)%Density                      = -1.0               ! [kg/m^3]
  BC_Input(1)%Mach                         = -1.0
  BC_Input(1)%VelocityX                    = 0.0                ! [m/s]
  BC_Input(1)%VelocityY                    = 0.0                ! [m/s]
  BC_Input(1)%VelocityZ                    = 0.0                ! [m/s]
  BC_Input(1)%Total_Temperature            = -1.0               ! [K]
  BC_Input(1)%Total_Pressure               = -1.0               ! [Pa]
  BC_Input(1)%Total_Density                = -1.0               ! [kg/m^3]
  BC_Input(1)%AoA_Alpha                    = 0.0                ! [deg]
  BC_Input(1)%AoA_Beta                     = 0.0                ! [deg]
  BC_Input(1)%Wall_Temp                    = -1.0               ! [K]
  BC_Input(1)%Rotation_Rate(1:3)           = [0.0,0.0,0.0]      ! [rpm]
  BC_Input(1)%Is_Specified_In_MRF          = F
  BC_Input(1)%Mass_Frac(1:7)               = [7*0.0]
  BC_Input(1)%Species_Name(1:7)            = [7*""]
  BC_Input(1)%Relax%Value                  = 1.0
  BC_Input(1)%Relax%Time                   = 0.0                ! non-dim
  BC_Input(1)%Relax%Iter                   = 0
  BC_Input(1)%SEM%Profile_File             = ""
  BC_Input(1)%SEM%Is_Cylindrical           = F
  BC_Input(1)%SEM%Length_Scale             = 0.0                ! [m]
  BC_Input(1)%SEM%Convective_Velocity      = 0.0                ! [m/s]
  BC_Input(1)%SEM%Profile_Height           = 0.0                ! [m]
  BC_Input(1)%SEM%Sub_In_BC_Grp            = 0
  BC_Input(1)%SEM%Sub_In_BC_Idx            = 0
  BC_Input(1)%Periodic%Matching_BC_Name    = ""
  BC_Input(1)%Periodic%Matching_BC_Idx     = 0
  BC_Input(1)%Periodic%RotationCenter(1:3) = [0.0,0.0,0.0]      ! grid units
  BC_Input(1)%Periodic%RotationAngle(1:3)  = [0.0,0.0,0.0]      ! [deg]
  BC_Input(1)%Periodic%RotationAxis(1:4)   = [0.0,0.0,0.0,0.0]
  BC_Input(1)%Periodic%Translation(1:3)    = [0.0,0.0,0.0]      ! grid units
/
\end{nmlblocklst}
\inpBC*[nmlBC::BCInput]{}

%% These nml_bc entries were used to test \texttt and hyperlinks
%% within the namelist block. To make the standard lstlisting match
%% the \texttt font, must use ',basewidth={0.50em,0.45em}'.
%% An extra space or two before the '=' is gobbled by Latex escape \@,
%% which might affect alignment of '='.
%%
%  !@\texttt{BC\_Input(1)\%Periodic\%Translation(1:3)}!@      = [0.0,0.0,0.0]
%  !@\BCinpBC[nmlBC::BCInput/Periodic/Translation][1]{Periodic\%Translation(1:3)}!@    = [0.0,0.0,0.0]

%==========================================================================

\begin{nmldescitem}
Walls_Are_Exact = T
\end{nmldescitem}
\inp*[nmlBC::WallsAreExact]{}[Walls\_Are\_Exact][]

\begin{nmldescbody}
Specifies whether the wall boundary conditions are set
directly using the exact flow conditions or
indirectly using the ghost state that provides
the exact flow conditions when averaged with the interior state.
\\
\begin{tabular}{ @{\qquad} r @{ = } p{0.75\linewidth} @{} }
  T & The wall boundary solution is the exact wall boundary condition,
      \forexample
      $\left[ u_{\textrm{wall}},
              v_{\textrm{wall}},
              w_{\textrm{wall}}
       \right] = 0$ for a no-slip wall. \\
  F & The average of the wall boundary solution
      and the interior solution on the boundary face
      recovers the exact wall boundary condition,
      \forexample
      $\left[ u_{\textrm{wall}},
              v_{\textrm{wall}},
              w_{\textrm{wall}}
       \right] =
      -\left[ u_{\textrm{wall}},
              v_{\textrm{wall}},
              w_{\textrm{wall}}
       \right]$ for a no-slip wall.
\end{tabular}

\NOTE
Hartmann and Leicht\cite{Hartmann_2013} report (Remark 6.13 on page 83) that
using exact wall boundary conditions seems to be more accurate on coarse grids
and for lower values of
\linebreak[4]% Prevent next word from entering right margin.
\inp[nmlNumerics::SolutionOrder]{Solution\_Order},
whereas using the averaging method is more stable.
%
%Details can be found in \linecite{Hartmann_2013}.
%{\small \url{http://elib.dlr.de/90967/1/hartmann_leicht_VKI_LS_2014-3.pdf}}

%\FIXME{
%\\
%This description could be worded better.\\
%Is this setting the external BC state or the actual face point values?\\
%For each case above, at what location is
%      $\left[ u_{\textrm{wall}},
%              v_{\textrm{wall}},
%              w_{\textrm{wall}}
%       \right]$
%specified? Is it a ghost cell?\\
%}
\end{nmldescbody}

%==========================================================================

\begin{nmldescitem}
Sub_Inflow_Method = 4
\end{nmldescitem}
\inp*[nmlBC::SubInflowMethod]{}[Sub\_Inflow\_Method][]

\begin{nmldescbody}
Specifies which subsonic inflow boundary condition algorithm to use.
\\
\begin{tabular}{ @{\qquad} r @{ = } p{0.85\linewidth} @{} }
  1 & Hold the total pressure and total temperature constant at the inflow,
      use the outgoing characteristic from the interior, and
      perform a Newton iteration to get
      the static temperature, speed of sound, and normal velocity magnitude
      for the exterior state. \\
  2 & Use the BC static density and velocity, and the interior static
      pressure. \\
  3 & Use the method specified in
      \linecite{Carlson_2011}. % NASA TM-2011-217181.
      The interior static pressure, total enthalpy, and
      outward propagating Riemann invariant are extrapolated
      from the interior to the boundary.
      %
      These values with the BC total pressure and total temperature
      are used to compute the BC speed of sound
      by solving a quadratic equation based on adiabatic and
      isentropic assumptions\slash relations.\\
  4 & Use the BC total pressure, total temperature, angle of attack,
      and the interior state to get the interior Mach number at the boundary.
\end{tabular}
\\
See \sref{SEC:BC:Inflow:Subsonic} for a detailed description of how each of
these methods is implemented.

\NOTE
It is recommended to use options 4, 3, or 1 (in that order).
\end{nmldescbody}

%==========================================================================

\begin{nmldescitem}
Use_BC_Conflicts = T
\end{nmldescitem}
\inp*[nmlBC::UseBCConflicts]{}[Use\_BC\_Conflicts][]

\begin{nmldescbody}
Logical flag for determining whether or not to overwrite boundary solutions
in adjacent cells if co-located boundary points have different boundary
conditions.
%
The boundary preference hierarchy can be found in \file{kind\_types\_mod.f90} in
the array \file{bc\_hierarchy}.

\NOTE
This is only relevant if using Legendre-Gauss-Lobatto points
for solution points
($\inp[nmlNumerics::LocSolutionPts]{Loc\_Solution\_Pts}=2$).
\end{nmldescbody}

%==========================================================================

\begin{nmldescitem}
BC_Input(1)%Name = ""
\end{nmldescitem}
\inpBC*[nmlBC::BCInput/Name]{}[Name][]

\begin{nmldescbody}
Character string that matches the name of a boundary group in the grid file.
%
See \sref{SEC:Files:Grid_Input:CGNS:PWBC} for using Pointwise
to set boundary conditions for CGNS files.
%
This input is the means of connecting boundary faces in the grid file with a
boundary condition provided in the input file to specify the boundary type and
flow conditions to use.

\NOTE
The character string is treated as case insensitive.
\end{nmldescbody}

%==========================================================================

\FINALFIX{\newpage}
\begin{nmldescitem}
BC_Input(1)%Set_BC_Value = -9999
\end{nmldescitem}
\inpBC*[nmlBC::BCInput/SetBCValue]{}[Set\_BC\_Value][]

\begin{nmldescbody}
This overrides the boundary condition type specified in the grid file,
using this boundary condition type for the boundary group given by
\BCinptext[1]{Name}.
%
Valid values for boundary condition types are:\\ %[0.25em]

\begingroup
\NewDocumentCommand{\bcgroup}{m}{\textbf{#1}}
%
%% Use threeparttable to add table footnote.
\begin{threeparttable}
\begin{tabular}{ @{} r @{ \quad } p{0.845\linewidth} @{} }
%% Use longtable to permit page breaks.
%\begin{longtable}{ @{} r @{ \quad } p{0.845\linewidth} @{} }
        & \bcgroup{Inflow boundary conditions} \\*
      1 & Generic inflow, subsonic\slash supersonic  \\*
      2 & Subsonic inflow \\*
      3 & Supersonic inflow \\*
      5 & SEM (unsteady turbulent) inflow%
          \tnote{\ensuremath{\ast}} %
          \\[0.25em]
        & \bcgroup{Outflow boundary conditions} \\*
     11 & Generic outflow, subsonic\slash supersonic \\*
     12 & Subsonic outflow \\*
     13 & Supersonic outflow \\[0.25em]
        & \bcgroup{Farfield boundary conditions} \\*
     21 & Characteristic inflow\slash outflow based on Riemann invariants,
          subsonic\slash supersonic%
          \tnote{\ensuremath{\ast}} \\*
     22 & Generic inflow\slash outflow, subsonic\slash supersonic \\*
     23 & Fixed to freestream\slash reference flow conditions \\*
     24 & Fixed flow conditions (equivalent to supersonic inflow) \\[0.25em]
        & \bcgroup{Wall boundary conditions} \\*
     50 & Slip wall \\*
     52 & Adiabatic no-slip wall \\*
     53 & Isothermal no-slip wall \\[0.25em]
        & \bcgroup{Numerical boundary conditions} \\*
     80 & Symmetry boundary condition (equivalent to slip wall) \\*
     81 & Periodic (translated \andor rotated) coupled boundaries \\[0.25em]
        & \bcgroup{Other boundary conditions} \\*
  -9999 & Do not override the boundary condition type
          from the grid file. \\*
  \bottomrule
%\end{longtable}
\end{tabular}
  \begin{tablenotes}
    \item[\ensuremath{\ast}] Not compatible with multi-component gas mixtures.
  \end{tablenotes}
\end{threeparttable}
\endgroup
\end{nmldescbody}

%==========================================================================

\begin{nmldescitem}
BC_Input(1)%Temperature = -1.0  ! [K]
\end{nmldescitem}
\inpBC*[nmlBC::BCInput/Temperature]{}[Temperature][]

\begin{nmldescbody}
Static temperature (units of \unitTemp{}).
\end{nmldescbody}

%==========================================================================

\begin{nmldescitem}
BC_Input(1)%Pressure = -1.0  ! [Pa]
\end{nmldescitem}
\inpBC*[nmlBC::BCInput/Pressure]{}[Pressure][]

\begin{nmldescbody}
Static pressure (units of \unitPres{}).
\end{nmldescbody}

%==========================================================================

\begin{nmldescitem}
BC_Input(1)%Density = -1.0  ! [kg/m^3]
\end{nmldescitem}
\inpBC*[nmlBC::BCInput/Density]{}[Density][]

\begin{nmldescbody}
Static density (units of \unitDens{}).
\end{nmldescbody}

%==========================================================================

\begin{nmldescitem}
BC_Input(1)%Mach = -1.0
\end{nmldescitem}
\inpBC*[nmlBC::BCInput/Mach]{}[Mach][]

\begin{nmldescbody}
Mach number (dimensionless).
\end{nmldescbody}

%==========================================================================

\begin{nmldescitem}
BC_Input(1)%VelocityX = 0.0  ! [m/s]
\end{nmldescitem}
\inpBC*[nmlBC::BCInput/VelocityX]{}[VelocityX][]

\begin{nmldescbody}
Velocity in the $x$ coordinate direction (units of \unitVel{}).
\end{nmldescbody}

%==========================================================================

\begin{nmldescitem}
BC_Input(1)%VelocityY = 0.0  ! [m/s]
\end{nmldescitem}
\inpBC*[nmlBC::BCInput/VelocityY]{}[VelocityY][]

\begin{nmldescbody}
Velocity in the $y$ coordinate direction (units of \unitVel{}).
\end{nmldescbody}

%==========================================================================

\begin{nmldescitem}
BC_Input(1)%VelocityZ = 0.0  ! [m/s]
\end{nmldescitem}
\inpBC*[nmlBC::BCInput/VelocityZ]{}[VelocityZ][]

\begin{nmldescbody}
Velocity in the $z$ coordinate direction (units of \unitVel{}).
\end{nmldescbody}

%==========================================================================

\begin{nmldescitem}
BC_Input(1)%Total_Temperature = -1.0  ! [K]
\end{nmldescitem}
\inpBC*[nmlBC::BCInput/TotalTemperature]{}[Total\_Temperature][]

\begin{nmldescbody}
Total temperature (units of \unitTemp{}).
\end{nmldescbody}

%==========================================================================

\begin{nmldescitem}
BC_Input(1)%Total_Pressure = -1.0  ! [Pa]
\end{nmldescitem}
\inpBC*[nmlBC::BCInput/TotalPressure]{}[Total\_Pressure][]

\begin{nmldescbody}
Total pressure (units of \unitPres).
\end{nmldescbody}

%==========================================================================

\begin{nmldescitem}
BC_Input(1)%Total_Density = -1.0  ! [kg/m^3]
\end{nmldescitem}
\inpBC*[nmlBC::BCInput/TotalDensity]{}[Total\_Density][]

\begin{nmldescbody}
Total density (units of \unitDens{}).
\end{nmldescbody}

%==========================================================================

\begin{nmldescitem}
BC_Input(1)%AoA_Alpha = 0.0  ! [deg]
BC_Input(1)%AoA_Beta  = 0.0  ! [deg]
\end{nmldescitem}
\inpBC*[nmlBC::BCInput/AoAAlpha]{}[AoA\_Alpha][]
\inpBC*[nmlBC::BCInput/AoABeta]{}[AoA\_Beta][]

\begin{nmldescbody}
Flow angles (in degrees):
$\alpha$ is the
\key[key:AngleOfAttack]{angle of attack}[flow angles!angle of attack]
measured in the $x$-$y$ plane, and
$\beta$  is the
\key[key:AngleOfSideslip]{angle of sideslip}[flow angles!angle of sideslip]
measured in the $x$-$z$ plane.


%
These flow angles are defined by
\eref{EQN:FlowAngleVelX,EQN:FlowAngleVelY,EQN:FlowAngleVelZ} and
\fref{FIG:FlowAngle} on \pref{\pageref{FIG:FlowAngle}}.
\end{nmldescbody}

%==========================================================================

\begin{nmldescitem}
BC_Input(1)%Wall_Temp = -1.0  ! [K]
\end{nmldescitem}
\inpBC*[nmlBC::BCInput/WallTemp]{}[Wall\_Temp][]

\begin{nmldescbody}
Temperature for isothermal wall (units of \unitTemp{}).
\end{nmldescbody}

%==========================================================================

\begin{nmldescitem}
BC_Input(1)%Rotation_Rate(1:3) = [0.0,0.0,0.0]  ! [rpm] @(x,y,z)-axes
\end{nmldescitem}
\inpBC*[nmlBC::BCInput/RotationRate]{}[Rotation\_Rate][]

\begin{nmldescbody}
This specifies the rotation rate vector of a moving boundary surface.
%
The components are defined with respect to the Cartesian axes,
and the magnitude is the rotation rate (in RPMs).
\end{nmldescbody}

%==========================================================================

\begin{nmldescitem}
BC_Input(1)%Is_Specified_In_MRF = F
\end{nmldescitem}
\inpBC*[nmlBC::BCInput/IsSpecifiedInMRF]{}[Is\_Specified\_In\_MRF][]

\begin{nmldescbody}
Setting this to true indicates that this boundary condition is specified in
the moving reference frame. If the boundary is moving\slash rotating at the
same speed as the moving reference frame, this works as a shortcut to apply
the same motion to this boundary.
\end{nmldescbody}

%==========================================================================

\begin{nmldescitem}
BC_Input(1)%Mass_Frac(1:7) = [7*0.0]
\end{nmldescitem}
\inpBC*[nmlBC::BCInput/MassFrac]{}[Mass\_Frac][]

\begin{nmldescbody}
\raggedright % Avoid next word spilling in to right margin.
\BCinptext{Mass\_Frac(n)} specifies the mass fraction
of the species given in \BCinptext{Species\_Name(n)}.
\end{nmldescbody}

%==========================================================================

\begin{nmldescitem}
BC_Input(1)%Species_Name(1:7) = [7*""]
\end{nmldescitem}
\inpBC*[nmlBC::BCInput/SpeciesName]{}[Species\_Name][]

\begin{nmldescbody}
Specifies the chemical name\slash formula for each of the species.
\end{nmldescbody}

%==========================================================================

\subsubsection{Boundary Condition Relaxation Inputs}
\begin{nmldescitem}
BC_Input(1)%Relax%Value = 1.0
BC_Input(1)%Relax%Time  = 0.0  ! non-dim, t[s] * (a[m/s] / L=1.0[m])
BC_Input(1)%Relax%Iter  = 0
\end{nmldescitem}
\inpBC*[nmlBC::BCInput/Relax/Value]{}[Relax\%Value][]
\inpBC*[nmlBC::BCInput/Relax/Time]{}[Relax\%Time][]
\inpBC*[nmlBC::BCInput/Relax/Iter]{}[Relax\%Iter][]

\begin{nmldescbody}
Typically the volume solution is initialized from the
reference conditions specified in the
\\% Prevent next word from entering right margin.
\inp[nmlReference]{nml\_reference} namelist.
%
If boundary conditions are immediately applied, the sudden change
can introduce numerical instabilities.
%
These relaxation parameters provide a way to gradually apply
each boundary condition.

\BCinptext{Relax\%Value} specifies
the relaxation value (between 0.0 and 1.0) to limit the difference
between the interior and exterior states.
%
The default value of 1.0 does not impose any limitation.

There are two means for specifying when to end relaxation:
one based on solution time, the other based on solution iterations.
%
Relaxation will end when either criteria is met.

If $\BCinptext{Relax\%Time}>0.0$, it specifies
the total accumulated solution time
(nondimensionalized by the reference length of 1.0 meter and speed of sound)
to end relaxation.

If $\BCinptext{Relax\%Iter}>0$, it specifies
the total accumulated time step iteration to end relaxation.
\end{nmldescbody}


%==========================================================================

\subsubsection{Boundary Condition SEM Inputs}
\begin{nmldescitem}
BC_Input(1)%SEM%Profile_File = ""  ! 150 char max
\end{nmldescitem}
\inpBC*[nmlBC::BCInput/SEM/ProfileFile]{}[SEM\%Profile\_File][]

\begin{nmldescbody}
This gives the path to the file containing the flow profile used for the
synthetic eddy method (SEM) boundary condition.
%
See \sref{SEC:SEM:usage:profile} for a detailed description of the
required file format and flow variables needed for the SEM profile.
\end{nmldescbody}

%==========================================================================

\begin{nmldescitem}
BC_Input(1)%SEM%Is_Cylindrical = F
\end{nmldescitem}
\inpBC*[nmlBC::BCInput/SEM/IsCylindrical]{}[SEM\%Is\_Cylindrical][]

\begin{nmldescbody}
This is used to declare if the synthetic eddy method (SEM) box is cylindrical.
%
This solely affects how the eddies are randomly generated to achieve a
truly random distribution in the radial, $r$, direction.
%
For a cylindrical SEM box, randomly generating
  $r\ (r_{\textrm{min}}\le r \le r_{\textrm{max}})$ and
  $\theta\ (\theta_{\textrm{min}}\le \theta \le \theta_{\textrm{max}})$
coordinate values to place a new eddy will result in
more eddies near the symmetry axis and
less eddies further away.
%
In order to generate a truly uniform random distribution,
we need to use the equation
  $r=\sqrt{\mathtt{Random}()\left[ r^2_{\textrm{max}}
                                 - r^2_{\textrm{min}}
                            \right] + r^2_{\textrm{min}}}$
where $\mathtt{Random}()$ is a function that generates a random number
on the unit interval.
%
This applies an area weighting to the radial distribution and
results in a uniform distribution of eddies in the volume.

\end{nmldescbody}

%==========================================================================

\begin{nmldescitem}
BC_Input(1)%SEM%Length_Scale = 0.0  ! [m]
\end{nmldescitem}
\inpBC*[nmlBC::BCInput/SEM/LengthScale]{}[SEM\%Length\_Scale][]

\begin{nmldescbody}
This gives the length scale (in meters) of the ``zone of influence'' of each
synthetic eddy.
%
Recommended values are between $(1/3.5)\times$ and $(1/5)\times$ the
height of the boundary layer profile from the SEM profile file.
%
See \sref{SEC:SEM:usage:recommendations}
for a more detailed description of this input and how it relates to
\inpBC[nmlBC::BCInput/SEM/ConvectiveVelocity]{SEM\%Convective\_Velocity}.

\NOTE
Previous testing seems to suggest that a slower convective velocity
should be used when using a smaller length scale, and a higher convective
velocity should be used when using a larger length scale.
\end{nmldescbody}

%==========================================================================

\FINALFIX{\newpage}
\begin{nmldescitem}
BC_Input(1)%SEM%Convective_Velocity = 0.0  ! [m/s]
\end{nmldescitem}
\inpBC*[nmlBC::BCInput/SEM/ConvectiveVelocity]{}[SEM\%Convective\_Velocity][]

\begin{nmldescbody}
This gives the velocity (in m\slash s) at which the synthetic eddies are
convected through the SEM box.
%
Recommended values are between $0.25\times$ and $1\times$ the
freestream velocity for this BC.
%
See \sref{SEC:SEM:usage:recommendations}
for a more detailed description of this input and how it relates to
\\ % Prevent next word from entering right margin.
\inpBC[nmlBC::BCInput/SEM/LengthScale]{SEM\%Length\_Scale}.

\NOTE
Previous testing seems to suggest that a slower convective velocity
should be used when using a smaller length scale, and a higher convective
velocity should be used when using a larger length scale.
\end{nmldescbody}

%==========================================================================

\begin{nmldescitem}
BC_Input(1)%SEM%Profile_Height = 0.0  ! [m]
\end{nmldescitem}
\inpBC*[nmlBC::BCInput/SEM/ProfileHeight]{}[SEM\%Profile\_Height][]

\begin{nmldescbody}
\NOTE
\emph{This variable is not an input.}
%
It will be calculated based on the profile data read from
\inpBC[nmlBC::BCInput/SEM/ProfileFile]{SEM\%Profile\_File},
and will be reported with the boundary condition information
written to the display output.
%
This output can be used to verify that the profile data was
properly interpreted.

\end{nmldescbody}

%==========================================================================

\begin{nmldescitem}
BC_Input(1)%SEM%Sub_In_BC_Grp = 0
BC_Input(1)%SEM%Sub_In_BC_Idx = 0
\end{nmldescitem}
\inpBC*[nmlBC::BCInput/SEM/SubInBCGrp]{}[SEM\%Sub\_In\_BC\_Grp][]
\inpBC*[nmlBC::BCInput/SEM/SubInBCIdx]{}[SEM\%Sub\_In\_BC\_Idx][]

\begin{nmldescbody}
\NOTE
\emph{Theses variables are not inputs.}
%
They will be determined after setting up the SEM boundary condition
and will be reported with the boundary condition information
written to the display output.
%
This output can be used to verify that the profile data was
properly interpreted.

If the boundary layer profile specified in the
\inpBC[nmlBC::BCInput/SEM/ProfileFile]{SEM\%Profile\_File}
does not fully span the boundary domain, then
a new inflow BC will automatically be created for the boundary points
that are outside the SEM profile.
% 
\inptext{Sub\_In\_BC\_Grp} identifies the \emph{index of the BC group}
to which the new inflow BC created for this SEM belongs.
%
\inptext{Sub\_In\_BC\_Idx} identifies the \emph{index for the new inflow BC}
created for this SEM.

\NOTE
Inflow boundaries are automatically split in this way because
the synthetic eddy method of imposing inflow turublence
actually overprescribes flow conditions on subsonic boundaries.
%
Separating points outside the boundary layer profile into a
standard inflow boundary allows those points to properly communicate
information upstream as needed.
%
Therefore, we recommend that the SEM profile file not extend much
beyond the edge of the boundary layer.
%
See \sref{SEC:SEM:usage:subsonic} for a more detailed explanation.

\end{nmldescbody}

%==========================================================================

\subsubsection{Boundary Condition Periodicity Inputs}
\begin{nmldescitem}
BC_Input(1)%Periodic%Matching_BC_Name = ""
\end{nmldescitem}
\inpBC*[nmlBC::BCInput/Periodic/MatchingBCName]{}[Periodic\%Matching\_BC\_Name][]

\begin{nmldescbody}
This specifies the name of the other boundary condition that is a matching
periodic pair with this boundary condition.
%
If BC's $i$ and $j$ are a periodic pair, then
the following two relations need to be true:
\\[0.5em] % A little extra separation before tabular.
\begin{tabular}{ @{\qquad} l } % @{ \quad } p{0.85\linewidth} @{} }
  \BCinptext[i]{Periodic\%Matching\_BC\_Name} =
  \BCinptext[j]{Name}, and
\\
  \BCinptext[j]{Periodic\%Matching\_BC\_Name} =
  \BCinptext[i]{Name}.
\end{tabular}

\vspace{0.5em} % A little extra separation after tabular.
\NOTE
This and
\inpBC[nmlBC::BCInput/Periodic/MatchingBCIdx]{Periodic\%Matching\_BC\_Idx}
are redundant, but together they serve as a means of error checking the
inputs for periodic pairs.
\end{nmldescbody}

%==========================================================================

\FINALFIX{\newpage}
\begin{nmldescitem}
BC_Input(1)%Periodic%Matching_BC_Idx = 0
\end{nmldescitem}
\inpBC*[nmlBC::BCInput/Periodic/MatchingBCIdx]{}[Periodic\%Matching\_BC\_Idx][]

\begin{nmldescbody}
This specifies the index of the other boundary condition that is a matching
periodic pair with this boundary condition.
%
If BC's $i$ and $j$ are a periodic pair, then
the following two relations need to be true:
\\[0.5em] % A little extra separation before tabular.
\begin{tabular}{ @{\qquad} l } % @{ \quad } p{0.85\linewidth} @{} }
  \BCinptext[i]{Periodic\%Matching\_BC\_Idx} = j, and \\
  \BCinptext[j]{Periodic\%Matching\_BC\_Idx} = i.
\end{tabular}

\vspace{0.5em} % A little extra separation after tabular.
\NOTE
This and
\inpBC[nmlBC::BCInput/Periodic/MatchingBCName]{Periodic\%Matching\_BC\_Name}
are redundant, but together they serve as a means of error checking the
inputs for periodic pairs.
\end{nmldescbody}

%==========================================================================

\begin{nmldescitem}
BC_Input(1)%Periodic%RotationCenter(1:3) = [0.0,0.0,0.0]  ! [grid units]
\end{nmldescitem}
\inpBC*[nmlBC::BCInput/Periodic/RotationCenter]{}[Periodic\%RotationCenter][]

\begin{nmldescbody}
This array gives the coordinate origin
(in \key[key:GridUnits]{grid units}) used for
defining the rotation angles between rotated periodic boundary conditions.

\NOTE
This input should be the same for a matching pair of rotated periodic
boundary conditions.

\end{nmldescbody}

%==========================================================================

\begin{nmldescitem}
BC_Input(1)%Periodic%RotationAngle(1:3) = [0.0,0.0,0.0]  ! [deg]
\end{nmldescitem}
\inpBC*[nmlBC::BCInput/Periodic/RotationAngle]{}[Periodic\%RotationAngle][]

\begin{nmldescbody}
This array gives the rotation angles (in degrees) about each of the
coordinate axes from the current periodic boundary interface to the
connecting periodic boundary interface.
%
If rotating about more than one axis, the rotation is performed
first about the $x$-axis, then the $y$-axis, and finally the $z$-axis.

\NOTE
One may use either
\inptext{Periodic\%RotationAngle} or  % Short Form
%\BCinptext{Periodic\%RotationAngle} or  % Long Form
%\\ % Prevent next word from entering right margin. Needed with Long Form.
\inpBC[nmlBC::BCInput/Periodic/RotationAxis]{Periodic\%RotationAxis}
to define rotated periodic boundary conditions, but not both.
%
If both are specified, \inptext{RotationAxis} will be ignored and
\inptext{RotationAngle} will be used.

\NOTE If BC's $i$ and $j$ are a periodic pair, the following relation needs
to be true:
\\[0.5em] % A little extra separation before tabular.
\begin{tabular}{ @{\qquad} >{\small}l }
   \BCinptext[i]{Periodic\%RotationAngle(:)} = 
  -\BCinptext[j]{Periodic\%RotationAngle(:)}
\end{tabular}
\end{nmldescbody}

%==========================================================================

\begin{nmldescitem}
BC_Input(1)%Periodic%RotationAxis(1:4) = [0.0,0.0,0.0,0.0]
\end{nmldescitem}
\inpBC*[nmlBC::BCInput/Periodic/RotationAxis]{}[Periodic\%RotationAxis][]

\begin{nmldescbody}
The first 3 values of this array define a vector representing the axis of
rotation for this rotated periodic boundary condition, and
the last value of the array defines the rotation angle (in degrees) about
this axis.

\NOTE
The coordinate origin is used as the base of the vector defining the
axis of rotation, NOT
\inpBC[nmlBC::BCInput/Periodic/RotationCenter]{Periodic\%RotationCenter}.

\NOTE
The magnitude of the vector defining the axis of rotation is not used.
The sole purpose of this vector is to define the \emph{direction} of the
rotation axis.

\NOTE
One may use either
\inpBC[nmlBC::BCInput/Periodic/RotationAngle]{Periodic\%RotationAngle} or
%\\ % Prevent next word from entering right margin. Needed for Long Form.
%\BCinptext{Periodic\%RotationAxis} % Long Form
\inptext{Periodic\%RotationAxis} % Short Form
to define rotated periodic boundary conditions, but not both.
%
If both are specified, \inptext{RotationAxis} will be ignored and
\inptext{RotationAngle} will be used.

\NOTE
If BC's $i$ and $j$ are a periodic pair, one of the two following
sets of relations needs to be true:
\\[0.5em] % A little extra separation before tabular.
{\small
\begin{tabular}{ @{\qquad} l @{ = } l @{} l}
     \BCinptext[i]{Periodic\%RotationAxis(1:3)}
  & &\BCinptext[j]{Periodic\%RotationAxis(1:3)}
     \\
     \BCinptext[i]{Periodic\%RotationAxis(4)}
  &-&\BCinptext[j]{Periodic\%RotationAxis(4)}
\end{tabular}
}
\\[0.5em] % A little extra separation after tabular.
or
\\[0.5em] % A little extra separation before tabular.
{\small
\begin{tabular}{ @{\qquad} l @{ = } l @{} l}
     \BCinptext[i]{Periodic\%RotationAxis(1:3)}
  &-&\BCinptext[j]{Periodic\%RotationAxis(1:3)}
     \\
     \BCinptext[i]{Periodic\%RotationAxis(4)}
  & &\BCinptext[j]{Periodic\%RotationAxis(4)}
\end{tabular}
}
\\[0.5em] % A little extra separation after tabular.
In other words, the opposing periodic face must either negate the
axis direction or the rotation angle, but not both.
\end{nmldescbody}

%==========================================================================

\begin{nmldescitem}
BC_Input(1)%Periodic%Translation(1:3) = [0.0,0.0,0.0]  ! [grid units]
\end{nmldescitem}
\inpBC*[nmlBC::BCInput/Periodic/Translation]{}[Periodic\%Translation][]

\begin{nmldescbody}
This array gives the translation vector
(in \key[key:GridUnits]{grid units}) from the current
periodic boundary interface to the connecting periodic boundary interface.

\NOTE
If BC's $i$ and $j$ are a periodic pair, the following relation needs
to be true:
\\[0.5em] % A little extra separation before tabular.
\begin{tabular}{ @{\qquad} l }
     \BCinptext[i]{Periodic\%Translation(:)} = 
    -\BCinptext[j]{Periodic\%Translation(:)}
\end{tabular}
\\[0.5em] % A little extra separation after tabular.
In other words, the opposing periodic face must negate the translation
vector.
\end{nmldescbody}

%==========================================================================

